\section{Introduction}

Ce chapitre présente la partie réalisée du projet. Il décrit l'environnement de travail, les technologies
utilisées, l'organisation du dépôt et les principales fonctionnalités développées dans le backend, le
dashboard administrateur et l'application mobile. La dernière partie est consacrée aux tests, à la
configuration et aux difficultés rencontrées pendant le développement.

\section{Environnement matériel et logiciel}

Le développement de PharmacieConnect a nécessité plusieurs outils, car le projet contient trois parties :
un backend, une interface web d'administration et une application mobile.

\begin{table}[H]
    \centering
    \caption{Environnement de développement}
    \begin{tabular}{L{5cm}L{8cm}}
        \toprule
        \textbf{Élément} & \textbf{Description} \\
        \midrule
        Machine de développement & [À compléter] \\
        Système d'exploitation & [À compléter] \\
        Éditeur de code & [À compléter] \\
        Backend & Python, FastAPI, Uvicorn, SQLAlchemy, Pydantic, Pytest \\
        Administration web & Node.js, React, Vite, React Router, Axios, Tailwind CSS \\
        Application mobile & Expo, React Native, React Navigation, i18next \\
        Base de données & SQLite en développement local, PostgreSQL selon la configuration \\
        Outils de test & Pytest, tests API, tests fonctionnels manuels \\
        \bottomrule
    \end{tabular}
\end{table}

\section{Technologies utilisées}

Les technologies ont été choisies en fonction des besoins du projet. Le backend devait exposer des API
claires et sécurisées. L'administration devait permettre la gestion des données depuis un navigateur.
L'application mobile devait rester adaptée à un usage simple et rapide.

\begin{longtable}{L{4cm}L{10cm}}
    \caption{Technologies utilisées dans le projet}\\
    \toprule
    \textbf{Technologie} & \textbf{Utilisation} \\
    \midrule
    \endfirsthead
    \toprule
    \textbf{Technologie} & \textbf{Utilisation} \\
    \midrule
    \endhead
    FastAPI & Développement de l'API REST, endpoints publics et endpoints protégés. \\
    SQLAlchemy & Modélisation des entités et accès à la base de données. \\
    Pydantic & Validation des données entrantes et structuration des réponses. \\
    JWT & Authentification, jetons d'accès et jetons de rafraîchissement. \\
    React & Construction de l'interface d'administration. \\
    Vite & Outil de développement et de build pour le dashboard administrateur. \\
    Axios & Communication entre le dashboard administrateur et l'API. \\
    Tailwind CSS & Mise en forme de l'interface d'administration. \\
    Expo & Développement et exécution de l'application mobile. \\
    React Native & Création des écrans et composants mobiles. \\
    i18next & Gestion de l'internationalisation. \\
    Pytest & Tests automatisés du backend. \\
    \bottomrule
\end{longtable}

\section{Organisation du dépôt}

Le projet est organisé dans un même espace de travail. Cette organisation a facilité les tests entre les
trois parties, surtout pendant l'intégration entre l'API, le dashboard et l'application mobile.

\begin{table}[H]
    \centering
    \caption{Organisation générale du dépôt}
    \begin{tabular}{L{5.6cm}L{7.4cm}}
        \toprule
        \textbf{Dossier} & \textbf{Contenu} \\
        \midrule
        \path{backend_pharmacie/} & API FastAPI, modèles SQLAlchemy, services métier, migrations, sécurité et tests. \\
        \path{admin_pharmacie/} & Tableau de bord React/Vite destiné à l'administration. \\
        \begin{minipage}[t]{5.4cm}\ttfamily\small ouerkema-pharmacieconnect-\linebreak 4c94773cce7d/\end{minipage} & Application mobile Expo/React Native destinée aux utilisateurs finaux. \\
        \path{docs/} & Documentation technique liée à l'architecture, l'API, le backend, le web admin et le mobile. \\
        \bottomrule
    \end{tabular}
\end{table}

\section{Réalisation du backend}

Le backend est la partie centrale de PharmacieConnect. Il reçoit les requêtes des deux interfaces, applique
les règles métier et interagit avec la base de données. Il contient aussi les mécanismes d'authentification,
de validation, d'audit et d'import.

\subsection{Endpoints publics}

Les endpoints publics sont utilisés principalement par l'application mobile. Ils permettent de consulter
les informations sans passer par l'administration.

\begin{itemize}
    \item \texttt{GET /health} pour vérifier l'état du backend ;
    \item \texttt{GET /api/pharmacies} pour consulter les pharmacies avec pagination ;
    \item \texttt{GET /api/pharmacies/search} pour rechercher par texte ;
    \item \texttt{GET /api/pharmacies/nearby} pour rechercher les pharmacies proches d'une position GPS ;
    \item \texttt{GET /api/pharmacies/\{pharmacy\_id\}} pour consulter le détail d'une pharmacie ;
    \item \texttt{GET /api/gardes} pour consulter les gardes d'une date donnée ;
    \item \texttt{GET /api/medicines} pour rechercher dans le catalogue de médicaments ;
    \item \texttt{GET /api/medicines/\{code\_pct\}} pour consulter le détail d'un médicament.
\end{itemize}

Pour la recherche par proximité, l'API utilise les coordonnées GPS envoyées par le client. Elle calcule la
distance avec les pharmacies disponibles, puis retourne les résultats triés par distance.

\subsection{Authentification et sessions}

L'authentification est gérée côté backend afin de garder le contrôle sur les comptes et les autorisations.
Lorsqu'une connexion réussit, le serveur retourne un jeton d'accès et un jeton de rafraîchissement.

\begin{lstlisting}[caption={Structure simplifiée d'une réponse de connexion}]
{
  "access_token": "[jeton d'acces]",
  "refresh_token": "[jeton de rafraichissement]",
  "token_type": "bearer",
  "expires_in": 1800
}
\end{lstlisting}

Les routes protégées vérifient ensuite le jeton reçu et le rôle de l'utilisateur. Cette vérification est faite
côté serveur pour éviter qu'un simple contrôle d'interface ne soit contourné.

\subsection{Services métier}

La logique principale est placée dans des services. Ce choix évite de mettre trop de code dans les routes
et rend les traitements plus faciles à tester.

\begin{itemize}
    \item \texttt{AuthService} pour l'inscription, la connexion, la vérification email, le rafraîchissement et la déconnexion ;
    \item \texttt{PharmacyService} pour la gestion et l'import des pharmacies ;
    \item \texttt{GardeService} pour la gestion et l'import des plannings de garde ;
    \item \texttt{MedicineService} pour l'import, la recherche et la consultation des médicaments ;
    \item \texttt{AuditService} pour la journalisation des actions importantes ;
    \item \texttt{CacheService} pour l'utilisation du cache lorsque la configuration le permet.
\end{itemize}

\subsection{Imports CSV}

Les imports CSV ont été ajoutés pour éviter la saisie manuelle d'un grand nombre de lignes. Le backend
vérifie le fichier avant de l'enregistrer : extension, taille, colonnes obligatoires, formats des valeurs et
contraintes métier.

Lorsqu'une ligne est incorrecte, le système retourne un message d'erreur avec l'information nécessaire
pour corriger le fichier. Pour les médicaments, le code PCT permet de mettre à jour un enregistrement
existant au lieu de créer un doublon.

\subsection{Audit et analytics}

Les actions sensibles sont enregistrées dans les journaux d'audit. Cela concerne notamment les
connexions, les échecs de connexion, les inscriptions, les créations administratives et les imports. Les
données d'analytics sont utilisées pour alimenter le tableau de bord et donner une idée de l'activité du
système.

\section{Réalisation du dashboard administrateur}

Le dashboard administrateur permet de gérer les données du système depuis une interface web. Il est
destiné aux administrateurs et aux utilisateurs autorisés.

\subsection{Pages principales}

Les principales pages développées sont :

\begin{itemize}
    \item le tableau de bord pour les compteurs, indicateurs et activités récentes ;
    \item la page des pharmacies pour consulter, créer, modifier, supprimer et importer les pharmacies ;
    \item la page des gardes pour gérer les plannings de garde ;
    \item la page des médicaments pour importer et consulter le catalogue ;
    \item la carte pour visualiser les pharmacies géographiquement ;
    \item la page de management pour gérer les assistants et les autorisations ;
    \item la page des journaux d'audit pour consulter les actions enregistrées ;
    \item les pages de profil, paramètres, notifications et langues.
\end{itemize}

\subsection{Gestion des sessions côté administration}

Le dashboard utilise un client API centralisé. Ce client ajoute le jeton JWT aux requêtes protégées. Si le
jeton d'accès expire, il tente de renouveler la session. Si le renouvellement échoue, l'utilisateur est
redirigé vers la page de connexion.

\placeholderfigure
{Capture d'écran du tableau de bord administrateur}
{Capture à insérer : vue principale du dashboard avec les indicateurs et les accès rapides.}
{fig:screenshot-admin-dashboard}

\placeholderfigure
{Capture d'écran de la gestion des pharmacies}
{Capture à insérer : tableau des pharmacies avec actions de création, modification, suppression et import.}
{fig:screenshot-admin-pharmacies}

\section{Réalisation de l'application mobile}

L'application mobile est la partie utilisée par le citoyen. Elle regroupe les fonctionnalités de recherche,
de consultation des gardes, de carte, de médicaments et de paramètres.

\subsection{Recherche et détail des pharmacies}

L'écran d'accueil permet de rechercher une pharmacie par texte et de consulter les résultats. Le détail
d'une pharmacie affiche les informations disponibles, comme le nom, l'adresse, le téléphone, le
gouvernorat et les coordonnées lorsque celles-ci existent.

\subsection{Carte et géolocalisation}

L'application peut demander l'autorisation d'utiliser la position de l'utilisateur. Cette position est envoyée
au backend pour récupérer les pharmacies proches. La carte permet ensuite d'afficher les résultats sous
forme de marqueurs.

\subsection{Calendrier des gardes}

L'écran calendrier permet de choisir une date et de consulter les pharmacies de garde correspondantes.
Cette fonctionnalité répond au besoin principal du projet : trouver rapidement une pharmacie disponible
selon une date précise.

\subsection{Catalogue de médicaments}

L'écran des médicaments permet de rechercher par nom commercial, DCI ou code PCT. La fiche de détail
affiche les informations présentes dans la base, comme le prix public, le tarif de référence et la catégorie
de remboursement.

\subsection{Paramètres, langues et thème}

L'application prend en charge le français, l'anglais et l'arabe. Le choix de l'arabe active l'affichage RTL.
Les paramètres regroupent aussi le thème, les notifications, les favoris et certaines préférences
utilisateur.

\placeholderfigure
{Capture d'écran de l'accueil mobile}
{Capture à insérer : écran de recherche des pharmacies dans l'application mobile.}
{fig:screenshot-mobile-home}

\placeholderfigure
{Capture d'écran du calendrier des gardes}
{Capture à insérer : écran mobile de consultation des pharmacies de garde par date.}
{fig:screenshot-mobile-gardes}

\placeholderfigure
{Capture d'écran du catalogue de médicaments}
{Capture à insérer : écran mobile de recherche et consultation des médicaments.}
{fig:screenshot-mobile-medicines}

\section{Intégration API}

L'intégration API a relié les deux interfaces au backend. Le dashboard administrateur utilise Axios, tandis
que l'application mobile possède sa propre configuration d'appel API. Les deux clients consomment les
endpoints FastAPI selon leurs besoins.

\begin{table}[H]
    \centering
    \caption{Principales intégrations API}
    \begin{tabular}{L{4cm}L{5cm}L{4cm}}
        \toprule
        \textbf{Client} & \textbf{Endpoints utilisés} & \textbf{Objectif} \\
        \midrule
        Application mobile & Pharmacies, gardes, médicaments, authentification & Consultation et parcours citoyen \\
        Dashboard admin & Administration, imports, analytics, audit, authentification & Gestion opérationnelle \\
        Backend & Base de données, services métier, sécurité & Centralisation et validation \\
        \bottomrule
    \end{tabular}
\end{table}

\section{Déploiement et configuration}

La configuration du backend repose sur des variables d'environnement. Cela permet de séparer le code des
valeurs qui changent selon l'environnement, comme la base de données ou les clés de sécurité.

\begin{table}[H]
    \centering
    \caption{Variables d'environnement principales}
    \begin{tabular}{L{6.2cm}L{6.8cm}}
        \toprule
        \textbf{Variable} & \textbf{Rôle} \\
        \midrule
        \texttt{DATABASE\_URL} & Chaîne de connexion à la base de données. \\
        \texttt{SECRET\_KEY} & Clé utilisée pour signer les jetons JWT. \\
        \texttt{ALGORITHM} & Algorithme de signature des jetons. \\
        \texttt{ACCESS\_TOKEN\_EXPIRE\_MINUTES} & Durée de validité du jeton d'accès. \\
        \texttt{REFRESH\_TOKEN\_EXPIRE\_DAYS} & Durée de validité du jeton de rafraîchissement. \\
        \texttt{FRONTEND\_URL} & Origine autorisée pour l'interface web. \\
        \texttt{TRUSTED\_HOSTS} & Liste des hôtes autorisés. \\
        \texttt{ENABLE\_REDIS\_CACHE} & Activation ou désactivation du cache Redis. \\
        \bottomrule
    \end{tabular}
\end{table}

L'environnement exact de déploiement final est [À compléter]. Les commandes de lancement local sont
placées en annexe pour ne pas alourdir le corps du rapport.

\section{Tests et validation}

\subsection{Stratégie de test}

La validation a combiné des tests automatisés côté backend et des tests manuels sur les interfaces. Les
tests backend ont été utilisés pour vérifier les règles les plus sensibles, notamment l'authentification,
l'audit, les restrictions régionales, les analytics et les imports.

\begin{table}[H]
    \centering
    \caption{Stratégie de validation}
    \begin{tabular}{L{4cm}L{9cm}}
        \toprule
        \textbf{Type de test} & \textbf{Objectif} \\
        \midrule
        Tests unitaires backend & Vérifier une règle métier isolée dans un service. \\
        Tests API & Vérifier les endpoints, les statuts HTTP, les réponses et les droits d'accès. \\
        Tests d'import & Contrôler les fichiers CSV, les colonnes, les lignes invalides et les bilans d'import. \\
        Tests de sécurité & Vérifier les rôles, la vérification email, les jetons et les accès interdits. \\
        Tests manuels web & Valider les pages d'administration, formulaires, tableaux, cartes et imports. \\
        Tests manuels mobile & Valider la recherche, la carte, le calendrier, les langues, les paramètres et les médicaments. \\
        \bottomrule
    \end{tabular}
\end{table}

\subsection{Tests automatisés du backend}

Les tests backend sont écrits avec Pytest. Ils utilisent une base de test afin de ne pas modifier les données
de développement. Les fichiers de test identifiés couvrent notamment :

\begin{itemize}
    \item \texttt{test\_analytics\_dashboard.py} pour les indicateurs du tableau de bord ;
    \item \texttt{test\_assistant\_regions.py} pour les restrictions régionales des assistants ;
    \item \texttt{test\_audit\_logs.py} pour la traçabilité des connexions et des actions ;
    \item \texttt{test\_medicine\_upload.py} pour l'import et la recherche de médicaments.
\end{itemize}

\subsection{Tests API}

Les tests API vérifient les routes publiques et protégées. Ils contrôlent les paramètres requis, les limites
de pagination, les erreurs de validation, les accès interdits et la cohérence des réponses JSON.

\subsection{Validation fonctionnelle}

Les scénarios fonctionnels validés sont présentés dans le tableau suivant.

\begin{longtable}{L{4cm}L{5cm}L{5cm}}
    \caption{Scénarios fonctionnels testés}\\
    \toprule
    \textbf{Scénario} & \textbf{Étapes} & \textbf{Résultat attendu} \\
    \midrule
    \endfirsthead
    \toprule
    \textbf{Scénario} & \textbf{Étapes} & \textbf{Résultat attendu} \\
    \midrule
    \endhead
    Connexion administrateur & Saisir email et mot de passe, puis valider. & Accès au tableau de bord et stockage des jetons. \\
    Recherche mobile & Saisir un nom, une adresse ou un gouvernorat. & Affichage d'une liste de pharmacies pertinentes. \\
    Recherche par proximité & Autoriser la localisation et lancer la recherche. & Pharmacies triées par distance. \\
    Consultation des gardes & Choisir une date dans le calendrier. & Affichage des pharmacies de garde correspondantes. \\
    Import pharmacies & Charger un fichier CSV valide. & Données importées et bilan affiché. \\
    Import gardes & Charger un planning CSV. & Gardes consultables côté administration et côté mobile. \\
    Recherche médicament & Chercher par nom, DCI ou code PCT. & Résultats pertinents et détail accessible. \\
    Audit logs & Filtrer les actions enregistrées. & Résultats paginés cohérents. \\
    Changement de langue & Passer du français à l'arabe. & Interface traduite avec direction RTL appliquée. \\
    \bottomrule
\end{longtable}

\section{Difficultés rencontrées et solutions adoptées}

Plusieurs difficultés sont apparues pendant la réalisation. Elles ont surtout concerné la cohérence des
données, les droits d'accès et l'intégration entre les trois parties du projet.

\begin{longtable}{L{5cm}L{9cm}}
    \caption{Difficultés rencontrées et solutions adoptées}\\
    \toprule
    \textbf{Difficulté} & \textbf{Solution adoptée} \\
    \midrule
    \endfirsthead
    \toprule
    \textbf{Difficulté} & \textbf{Solution adoptée} \\
    \midrule
    \endhead
    Alignement entre backend, web et mobile & Centralisation des contrats API et séparation des services backend. \\
    Qualité variable des fichiers CSV & Validation des colonnes, contrôle des types, retour des erreurs ligne par ligne et comportement d'upsert pour les médicaments. \\
    Gestion des rôles et des assistants régionaux & Vérification des permissions côté backend et application du périmètre régional dans les services. \\
    Géolocalisation mobile & Demande explicite des permissions, contrôle de l'activation des services de position et recherche par coordonnées. \\
    Maintenabilité du backend & Séparation entre routes, services, modèles, schémas, sécurité et migrations. \\
    Internationalisation mobile & Utilisation de fichiers de traduction et prise en charge de la direction RTL pour l'arabe. \\
    \bottomrule
\end{longtable}

\section{Conclusion}

La réalisation a permis d'obtenir une plateforme fonctionnelle composée d'une API sécurisée, d'un
dashboard administrateur et d'une application mobile. Les tests effectués ont validé les principaux
parcours : authentification, recherche de pharmacies, consultation des gardes, import CSV, gestion des
rôles, audit et recherche de médicaments. Des améliorations restent possibles, notamment les tests
end-to-end, les tests de charge et la validation sur un volume de données plus important.
